\begin{prop}
    If $\phi : V \rightarrow V$ has a matrix $A$ with respect to a basis $\beta$. Then the following are equivalent:
    \begin{itemize}
      \item $\phi$ is diagonalisable
      \item $\phi_A$ is diagonalisable
      \item $A$ is diagonalisable
    \end{itemize}
\end{prop}
\begin{proof}
  $\phi_A$ and $\phi$ have the same eigenvalues and eigenvectors, so one being diagonalisable implies the other is too.

  \noindent If $A$ is diagonalisable, then
  \begin{align*}
    P^{-1} A P &= D \\
    A P &= P D
  \end{align*}
  and since $D$ is diagonal with diagonals $\lambda_1,...,\lambda_n$, $col_j(PD) = \lambda_j col_j(P)$. Thus
  \[col_j(AP) = \lambda_j col_j(P).\]
  Then $A (col_j(P)) = \lambda col_j(P)$, so each column of $P$ is an eigenvector. Since $P$ is invertible its columns are a basis thus $\phi_A$ is diagonalisable.
\end{proof}

\begin{remark}
    Many matrices are not diagonalisable over $\mathbb{R}$, but are diagonalisable over $\mathbb{C}$. This happens when the characteristic polynomial has no real roots.
\end{remark}

\subsection{Applications of Diagonalisation}

Diagonalisation can be extremely helpful for efficient calculation of powers of matrices.
\begin{example} $ $
    
Suppose $A$ is diagonalisable, with $P^{-1}AP = D$ being diagonal. Furthermore $P$ has columns $j$ equal to the $\lambda_j$ eigenvectors of $A$.
  Then $A = PDP^{-1}$, so for $k > 0$
  \begin{align*}
    A^k &= (PDP^{-1})(PDP^{-1})...(PDP^{-1}) \\
        &= PD^kP^{-1} \\
    \text{with } D^k = &\begin{pmatrix}
  \lambda_1^k & \dots & 0\\
  \vdots & \ddots & \vdots\\
  0 & \dots & \lambda_n^k\\
  \end{pmatrix}
  .\end{align*}
  This allows us to calculate $A^k$ explicitly with only $2$ matrix mutliplications.
\end{example}

\subsection{Multiplicity of Eigenvalues and Diagonalisation}

We want a necessary and sufficient condition for a matrix or operator to be diagonalisable.


If $\dim V = n$ and $\phi : V \rightarrow V$ is linear, then an eigenbasis needs $n$ linearly independent eigenvectors.

\begin{theorem}
    If $v_1,...,v_n$ are eigenvectors $\phi$ with distinct eigenvalues $\lambda_1,...,\lambda_n$ then $v_1,...,v_n$ are independent.
\end{theorem}
\begin{proof}
    By contradiction, suppose that $a_1 v_1 + ... + a_k v_k = 0$ is the smallest dependent list of eigenvectors. Then we have
    \begin{align*}
      \lambda_1 \sum_{ i = 1}^{k} a_i v_i = \lambda_1 \cdot 0 &= 0 \\
      \phi \left( \sum_{ i = 1}^{k} a_i v_i \right) = \phi(0) &= 0 
    .\end{align*}
    However $\phi(0)$ also equals
    \[ \phi \left( \sum_{ i = 1}^{k} a_i v_i \right) = \sum_{ i = 1}^{k} a_i \phi(v_i) = \sum_{ i = 1}^{n} a_i \lambda_i v_i.\]
    Then subtracting one from the other means the first term equals zero, meaning there exists a smaller dependence, a contradiction.
\end{proof}

\begin{remark}
    This can be generalized: if we conbine independent list from each eigenspace we get an independent list.
\end{remark}

\begin{corollary}
Let $n = \dim V$ and $\phi : V \rightarrow V$ be a linear operator. If 
\[\Delta_\phi(t) = (t - \lambda_1)...(t - \lambda_n)\]
with $\lambda_i$ distinct for all $i$, then $\phi$ is diagonalisable.
\end{corollary} 
\begin{proof}
  Because there are $n$ distinct eigenvalues, there are $n$ distinct eigenvectors and since they are independent, they form an eigenbasis.
\end{proof}


\begin{definition}
    Let $\lambda$ be an eigenvector of $\phi: V \rightarrow V$. Then
    \begin{itemize}
    \item the {\bf algebraic multiplicity} of $\lambda$ denoted $am(\lambda)$ is the largest $k$ such that $(\lambda - t)^k$ divides $\Delta_\phi(t)$.
    \item the {\bf geometric multiplicity} of $\lambda$ is $gm(\lambda) := \dim E_\phi(\lambda)$.
    \end{itemize}
\end{definition}

\begin{remark} $ $
  \begin{itemize}
  \item $\dim E_\phi(\lambda)$ is the maximal number of independent $\lambda-$eigenvectors.
  \item Suppose $\lambda$ is a root of $\Delta_\phi(t)$. Then $(t - \lambda)$ divides $\Delta_\phi(t)$ so
    \[am(\lambda) \geq 1.\]
    If $E_\phi(\lambda) \not = 0$ then 
    \[gm(\lambda) \geq 1.\]
  \end{itemize}    
\end{remark}

\begin{prop}
  For any $\lambda \in \mathbb{F}$ and any linear map $\phi$:
    \[am(\lambda) \geq gm(\lambda).\]
\end{prop}
\begin{proof}
  Let $v_1,...,v_k$ be a basis of $E_\phi(\lambda)$. Extend this list to be a basis of $V$. Then the matrix of $\phi$ from the standard basis to this basis is diagonal in the top corner, and so if $A$ is the matrix in the bottom corner, then 
  \[\Delta_\phi(t) = (\lambda - t)^t \det(A -It).\]
  Hence $(\lambda - t)$ divides $\Delta_\phi(t)$ at least $k$ times, so $am(\lambda) \geq gm(\lambda)$.
\end{proof}

\begin{theorem}
    Let $\phi : V \rightarrow V$ be a linear operator. Then $\phi$ is diagonalisable if and only if $\Delta_\phi(t)$ is a product of linear factors and $am(\lambda) = gm(\lambda)$ for all $\lambda$.
    \end{theorem}
\begin{proof} $ $ \\
  $(\implies)$ $\phi$ has an eigenbasis and is diagonalisable, so $\Delta_\phi(t)$ is clearly a product of linear factors. If $am(\lambda) > gm(\lambda)$ then there is a $\lambda-$eigenvector in this eigenbasis that can be written in terms of other $\lambda-$eigenvectors, thus contradicting that $\phi$ is diagonalisable, thus they are equal.
    $ $ \\
  $ $ \\
  \noindent $(\Leftarrow)$ As $\Delta_\phi(t)$ is a product of linear factors, we can write it as
  \[\Delta_\phi(t) = (\lambda_1 - t)^{k_1} ... (\lambda_m - t)^{k_m}\]
  and $k_1, ..., k_m$ sum to $n$. But $k_i = am(\lambda_i) = gm(\lambda_i)$. Thus $\sum_{ i = 1}^{m} gm(\lambda_i) = n$. Therefore we have a length $n$ list of eigenvectors and so if they are independent they form a basis and $\phi$ is diagonalisable.

  Suppose there is a dependence. Then we can group all $\lambda_i-$eigenvectors together as a single eigenvector $w_i$. Since there is now a sum of distinct eigenvectors there cannot be a dependence, so the dependence must be inside a $w_i$. However $w_i$ is formed from a basis so there also cannot be a dependence here. Thus $\phi$ is diagonalisable.

\end{proof}
